\documentclass[11pt]{article}
\usepackage[margin=1in]{geometry}
\usepackage{amsmath,amssymb,amsthm}
\usepackage[hidelinks]{hyperref}
\newcommand{\Tr}{\operatorname{Tr}}
\newcommand{\lmin}{\lambda_{\min}}
\title{A three-dimensional counterexample to Lin's\\
additive majorization conjecture\\
\large Complete solution of MathDB 332939}
\author{Principia Math in Progress}
\date{12 September 2026}
\begin{document}
\maketitle

\paragraph{Task 1. Resolve the all-power majorization conjecture.}
For positive semidefinite matrices $A,B$, let $|Z|=(Z^*Z)^{1/2}$ and
write eigenvalues in decreasing order. Lin's Conjecture 4.3 asks whether
\[
 \lambda(A^2+|BA|^p)\succ\lambda(A^2+|AB|^p)
 \qquad\text{for every }p>0.
\]
Here $\succ$ is ordinary majorization. The answer is negative, already
for $p=4$ and real positive definite matrices of order three.

Take
\[
 A=\operatorname{diag}(1,2,20),\qquad
 B=\frac15 M,\qquad
 M=\begin{pmatrix}6&-1&1\\-1&2&1\\1&1&1\end{pmatrix}.
\]
The leading principal minors of $M$ are $6,11,1$, so $B$ is positive
definite by Sylvester's criterion. The same is immediate for $A$.
Set
\[
 L=A^2+(AB^2A)^2=A^2+|BA|^4,\qquad
 R=A^2+(BA^2B)^2=A^2+|AB|^4.
\]
We shall certify
\begin{equation}\label{separation}
 \lmin(R)\leq\frac{7027}{1250}
 <\frac{142}{25}<\lmin(L).
\end{equation}
Since $L$ and $R$ have equal trace, this contradicts
$\lambda(L)\succ\lambda(R)$: the inequality for the two largest
eigenvalues would require $\lmin(L)\leq\lmin(R)$.

\paragraph{An exact lower bound for \(\lmin(L)\).}
Direct multiplication gives
\[
 X:=AM^2A=
 \begin{pmatrix}38&-14&120\\-14&24&80\\120&80&1200\end{pmatrix}.
\]
Consequently
\[
 T:=625L-3550I
 =625A^2+X^2-3550I
 =\begin{pmatrix}
 13115&8732&147440\\
 8732&6122&96240\\
 147440&96240&1707250
 \end{pmatrix}.
\]
The three leading principal minors of $T$ are exactly
\[
 13115,\qquad 4042206,\qquad 152317908700.
\]
They are positive, so $T$ is positive definite. Hence
$\lmin(L)>3550/625=142/25$.

\paragraph{An exact upper bound for \(\lmin(R)\).}
Put $v=(1,-1,0)^{\mathsf T}$ and $Y=MA^2M$. Then
\[
 Yv=\begin{pmatrix}54\\-31\\-5\end{pmatrix},\qquad
 625R=625A^2+Y^2.
\]
The Rayleigh quotient therefore gives
\[
 \lmin(R)\leq
 \frac{v^{\mathsf T}Rv}{v^{\mathsf T}v}
 =\frac{625\cdot5+54^2+31^2+5^2}{625\cdot2}
 =\frac{7027}{1250}.
\]
The separation is strict because
$142/25-7027/1250=73/1250>0$.
Finally, cyclicity of the trace gives $\Tr X^2=\Tr Y^2$;
in fact $\Tr(625L)=\Tr(625R)=1737137$.
Thus \eqref{separation} proves the failure of the required majorization.
\hfill$\square$

\paragraph{Relation to the cited even-power result.}
The earlier Corollary 2.3 of Lin and Wolkowicz proves an inequality
involving the words
$(AB)^k(BA)^k$ and $(BA)^k(AB)^k$.
For $k>1$ these are different from $|BA|^{2k}$ and $|AB|^{2k}$.
Our counterexample concerns the latter expressions, precisely as in
Conjecture 4.3. Ghabries's 2022 thesis \cite{G}, Remark 2.5, also records this
additive majorization question after proving its determinantal consequence.

\paragraph{Verification.}
The accompanying script checks all displayed products, positivity
certificates, trace identities and Rayleigh bounds over the rationals.
All thirteen checks pass. The proof uses no numerical eigenvalues.

\paragraph{Task 2. Determine the valid exponent range.}
Identify all $p>0$ for which the majorization statement holds for every
pair of positive semidefinite matrices.

\paragraph{Task 3. Construct systematic counterexample families.}
Find parameterized positive definite pairs for which the least-eigenvalue
inequality reverses, with explicit rational separation bounds.

\paragraph{Task 4. Identify structural sufficient conditions.}
Determine conditions beyond commutativity on $A$ and $B$ that guarantee
the majorization for every positive exponent.

\paragraph{Task 5. Determine surviving spectral inequalities.}
Characterize which convex spectral functionals retain the asserted
comparison even when full additive majorization fails.

\begin{thebibliography}{9}
\bibitem{Lin}
M. Lin, \emph{On a determinantal inequality arising from diffusion tensor
imaging}, Communications in Contemporary Mathematics \textbf{19} (2017),
1650044. Conjecture 4.3 in
\href{https://arxiv.org/abs/1604.04141}{arXiv:1604.04141}.
\bibitem{LW}
M. Lin and H. Wolkowicz, \emph{An eigenvalue majorization inequality for
positive semidefinite block matrices}, Linear and Multilinear Algebra
\textbf{60} (2012), 1365--1368. Corollary 2.3 in the
\href{https://www.math.uwaterloo.ca/~hwolkowi/henry/reports/finalLAMAdec18.pdf}{author's final manuscript}.
\bibitem{G}
M. M. Ghabries, \emph{Contributions to Matrix Inequalities and Some
Applications}, doctoral thesis, Universit\'e d'Angers, 2022,
Remark 2.5, p.~50,
\href{https://www.researchgate.net/publication/361793582_Contributions_to_Matrix_Inequalities_and_Some_Applications}{author-uploaded full text}.
\end{thebibliography}
\end{document}
